\chapter{Riflessione conclusiva}

\section{Soddisfacimento degli obiettivi}
\label{sec:soddisfacimento-obiettivi}

% In questa sezione verrà effettuata una valutazione oggettiva del grado di raggiungimento degli obiettivi definiti all'inizio del progetto di \textit{stage}. Per ciascun obiettivo verrà indicato se sia stato completamente raggiunto, parzialmente soddisfatto oppure non realizzato, motivando tali valutazioni sulla base di evidenze concrete. L'analisi farà riferimento ai risultati presentati nel Capitolo~3, considerando le funzionalità implementate, le piattaforme supportate, i requisiti di certificazione soddisfatti e le problematiche risolte durante il periodo di \textit{stage}.
 
Valuto in questa sezione il grado di raggiungimento degli obiettivi aziendali e personali definiti nel Capitolo~2 (Tabelle~\ref{tab:obiettivi_migrazione}, \ref{tab:obiettivi_switch} e \ref{tab:obiettivi_personali}), sulla base delle evidenze presentate nel Capitolo~3. Per ciascun obiettivo indico se sia stato completamente raggiunto, parzialmente raggiunto o non raggiunto; i dati quantitativi sottostanti sono già riassunti in Tabella~\ref{tab:confronto-obiettivi-risultati} (\S\ref{sec:certificazione-switch}) e non li ripeto qui.
 
\subsection{Obiettivi aziendali}
 
\begin{table}[h]
\centering
\begin{tabular}{|p{5cm}|p{3cm}|p{4cm}|}
\hline
\textbf{Obiettivo} & \textbf{Valutazione} & \textbf{Riferimento} \\
\hline
Migrazione Unity 2017 $\rightarrow$ Unity 6 & Raggiunto & \S3.2.1 \\
\hline
Stabilizzazione \textit{build} & Raggiunto & \S3.2.1 \\
\hline
Eliminazione \gls{OBB} \textit{mobile} & Raggiunto & \S3.2.1 \\
\hline
Compatibilità Steam & Raggiunto & \S3.2.1 \\
\hline
Navigazione \textit{controller} & Raggiunto & \S3.2.2 \\
\hline
Certificazione Nintendo Switch & Parzialmente raggiunto entro la durata del mio \textit{stage} & \S3.2.3 \\
\hline
\end{tabular}
\caption{Valutazione degli obiettivi aziendali}
\label{tab:valutazione-obiettivi-aziendali}
\end{table}
 
I primi cinque obiettivi non richiedono ulteriore commento: le evidenze riportate nei rispettivi paragrafi del Capitolo~3 ne confermano il pieno raggiungimento, senza eccezioni o compromessi rilevanti.
 
La certificazione Nintendo Switch merita invece una precisazione. Dei sei \textit{test} obbligatori (Tabella~\ref{tab:test-nintendo-salvataggi}), cinque sono stati superati durante le otto settimane del mio \textit{stage}, incluso il più impegnativo — il \textit{test} sul \textit{Write Access Count} — risolto tramite le tre ottimizzazioni descritte in \S\ref{sec:certificazione-switch}. Il sesto \textit{test}, relativo alla continuità del salvataggio al variare dell'\textit{account} Nintendo collegato, è stato completato la settimana successiva alla conclusione del mio \textit{stage} dal mio collega. Considerando l'intero arco del progetto, l'obiettivo aziendale risulta quindi pienamente soddisfatto; entro i tempi del mio \textit{stage} individuale, il grado di completamento è di 5 \textit{test} su 6.
 
\subsection{Obiettivi personali}
 
\begin{table}[h]
\centering
\begin{tabular}{|p{6cm}|p{3cm}|p{3cm}|}
\hline
\textbf{Obiettivo personale} & \textbf{Valutazione} & \textbf{Riferimento} \\
\hline
Acquisire competenza sulla gestione dell'\textit{input} da \textit{controller} & Raggiunto & \S3.2.2 \\
\hline
Comprendere il processo di certificazione per \textit{console} & Raggiunto & \S3.2.3 \\
\hline
Gestire la migrazione di una \textit{codebase} esistente & Raggiunto & \S3.2.1 \\
\hline
\end{tabular}
\caption{Valutazione degli obiettivi personali}
\label{tab:valutazione-obiettivi-personali}
\end{table}
 
Anche per l'obiettivo di comprendere il processo di certificazione vale la stessa precisazione fatta per l'obiettivo aziendale corrispondente: la comprensione del processo nel suo complesso, inclusa l'analisi e la risoluzione del problema di \textit{Write Access Count}, è avvenuta interamente durante il mio \textit{stage}; il caso specifico del sesto \textit{test} è stato invece compreso e risolto dal mio collega successivamente, e non rientra quindi nella mia esperienza diretta.

\section{Crescita professionale}

Lo \textit{stage} ha rappresentato un'importante occasione di crescita sia dal punto di vista tecnico sia da quello organizzativo. In questa sezione verranno descritte le competenze acquisite nel corso dell'esperienza, tra cui la padronanza del sistema di gestione degli \textit{input} di \textit{Unity} e delle relative peculiarità sulle piattaforme \textit{console}, la comprensione del processo di certificazione richiesto da un \textit{platform holder} quale \textit{Nintendo}, l'esperienza maturata nella gestione di una \textit{codebase legacy} caratterizzata dalla presenza di debito tecnico e lo sviluppo della capacità di lavorare in collaborazione con un altro sviluppatore su un progetto condiviso. Verrà inoltre distinta la formazione di nuove competenze dal consolidamento di conoscenze già possedute prima dell'inizio dello \textit{stage}.

\section{Distanza tra corso di studi e mondo del lavoro}

L'ultima sezione propone una riflessione sul rapporto tra il percorso universitario e le competenze richieste dall'attività svolta in azienda. Verranno analizzati gli aspetti per i quali la preparazione accademica si è dimostrata adeguata, come la conoscenza delle strutture dati, l'impiego di sistemi di versionamento in un contesto collaborativo, nonché la capacità di comprendere e analizzare codice sviluppato da altri programmatori. Parallelamente saranno evidenziate le principali competenze che hanno richiesto un apprendimento direttamente sul campo, tra cui l'utilizzo professionale di un motore grafico come \textit{Unity}, la gestione del ciclo di sviluppo di un videogioco destinato a piattaforme commerciali e la comprensione delle procedure di certificazione previste dai produttori di \textit{console}. La riflessione conclusiva metterà in evidenza come l'esperienza di stage abbia rappresentato un efficace punto di incontro tra la preparazione teorica acquisita durante il corso di laurea e le competenze pratiche richieste nel mondo dello sviluppo \textit{software} professionale.